% =====================================================================
% Paper 04 — Pillar-4 atomic theorem: SU(5) x U(1) emergence from BCC
%   topological defects on Sigma_0 = P^1 x P^1, with Cartan forcing,
%   triple verification, and F-GAP4 verdict shell. (Rev 3, 2026-05-03)
% Status: [DRAFT] [NEEDS_UPDATE] — Wave 4 gauge sector
% AUDIT-FLAG: AUDIT-2026-05-02-Wave7-Aux-Epoch-Overclaim
%   Rev 1 (2026-05-02): downgraded to "historical CP^2 auxiliary note"
%     per first-pass audit, but body theorem and F-GAP4 wording remained
%     incompatible with abstract.
%   Rev 2 (2026-05-02): partial fix; still left "fully rigorous" /
%     "5% precision" / "completes the proof" wording in body.
%   Rev 3 (2026-05-03): COMPLETE REWRITE with proper canonical content
%     synthesised from Math268+269 (original Pillar 4 atomic promotion),
%     Math276 (external defense programme, 5 attacks DISCHARGED/MITIGATED),
%     Math277+278 (third-order audits Lemma A/B, OUTCOME A both),
%     Math279 (final consolidation, T6 PROVED CONDITIONAL with explicit
%     hypothesis set H_atomic = {H1.1, H2.1-H2.7, H3.1, H_task}),
%     Math305 (explicit Sigma_0 = P^1 x P^1 4-chart Cech atlas + lifted
%     transition functions + Chern data), Math288 (Sigma_0 realization
%     readiness + defect-mass formula derivation + Bogomolov-type
%     eigenvalue bound), Math300+Math311 (F-GAP4-DEFECT-MASS verdict
%     shell: lambda_min(Box_E) -> mu_defect with PASS/FAIL/DEFER on
%     window [10^13, 10^17] GeV + solver/grid certificates, NOT 5%
%     precision prediction), Math229+275 (Cartan forcing first-principles
%     re-derivation), Math270 (cross-base coherence / topological
%     certificate transfer theorem), Math302 (proof-rigor vs
%     realisation-gate carve-out).
% Compile: pdflatex Paper-04.tex (×2 for refs)
% Canonical archive: Math162, 164, 170, 174, 191, 192, 229, 238, 240,
%   245, 248, 263, 266, 267, 268, 269, 270, 271, 272, 273, 274, 275,
%   276, 277, 278, 279, 280, 281, 286, 288, 300, 302, 305, 311,
%   Math314-AddD (audit)
% =====================================================================
% --- TECT canonical preamble (see _shared/tect-preamble.tex) ----------
% =====================================================================
% TECT canonical LaTeX preamble (revtex4-2 / single-column / theorem-safe)
% =====================================================================
% Maintainer: Jusang Lee (jtkor@outlook.com)
% Binding from: 2026-05-06
% Purpose: a single source-of-truth preamble for every TECT paper
% (Paper-00 .. Paper-10 + Auxiliary notes). Designed to (a) eliminate
% font-corruption on pdflatex (T1 fontenc + lmodern), (b) make
% theorem-heavy mathematical-physics content render cleanly in a
% single-column layout (PRD class option, NOT PRL two-column), (c)
% provide consistent theorem numbering across all TECT papers via
% amsthm with a shared counter set, (d) make hyperref + cleveref
% cross-references resolve cleanly with the standard two-pass
% pdflatex compile.
%
% Usage in each Paper-NN.tex:
%   % =====================================================================
% TECT canonical LaTeX preamble (revtex4-2 / single-column / theorem-safe)
% =====================================================================
% Maintainer: Jusang Lee (jtkor@outlook.com)
% Binding from: 2026-05-06
% Purpose: a single source-of-truth preamble for every TECT paper
% (Paper-00 .. Paper-10 + Auxiliary notes). Designed to (a) eliminate
% font-corruption on pdflatex (T1 fontenc + lmodern), (b) make
% theorem-heavy mathematical-physics content render cleanly in a
% single-column layout (PRD class option, NOT PRL two-column), (c)
% provide consistent theorem numbering across all TECT papers via
% amsthm with a shared counter set, (d) make hyperref + cleveref
% cross-references resolve cleanly with the standard two-pass
% pdflatex compile.
%
% Usage in each Paper-NN.tex:
%   % =====================================================================
% TECT canonical LaTeX preamble (revtex4-2 / single-column / theorem-safe)
% =====================================================================
% Maintainer: Jusang Lee (jtkor@outlook.com)
% Binding from: 2026-05-06
% Purpose: a single source-of-truth preamble for every TECT paper
% (Paper-00 .. Paper-10 + Auxiliary notes). Designed to (a) eliminate
% font-corruption on pdflatex (T1 fontenc + lmodern), (b) make
% theorem-heavy mathematical-physics content render cleanly in a
% single-column layout (PRD class option, NOT PRL two-column), (c)
% provide consistent theorem numbering across all TECT papers via
% amsthm with a shared counter set, (d) make hyperref + cleveref
% cross-references resolve cleanly with the standard two-pass
% pdflatex compile.
%
% Usage in each Paper-NN.tex:
%   \input{../_shared/tect-preamble.tex}
%   \begin{document}
%   ...
%   \end{document}
%
% Build (every paper): `pdflatex Paper-NN && pdflatex Paper-NN`
% (the second pass resolves \ref{} / \cref{} forward references).
% A bash/PowerShell wrapper that automates this is at
% Docs/papers/papers/_shared/build-paper.sh / build-paper.ps1.
% =====================================================================

% ---- Document class --------------------------------------------------
% PRD class option of revtex4-2 with [reprint,onecolumn] gives the same
% APS heading / by-line / abstract style as PRL but in a wider single
% column that handles long display equations and theorem environments
% without column-break artefacts. nofootinbib keeps citations out of
% footnotes. The 11pt option restores standard font metrics; with T1
% fontenc + lmodern this gives non-bitmap glyphs.
\documentclass[%
  aps,prd,reprint,onecolumn,nofootinbib,11pt,%
  amsmath,amssymb,floatfix,longbibliography%
]{revtex4-2}

% ---- Encoding & fonts ------------------------------------------------
\usepackage[T1]{fontenc}        % proper 8-bit font encoding (no font corruption)
\usepackage[utf8]{inputenc}     % allow UTF-8 source
\usepackage{lmodern}            % scalable Latin Modern (replaces bitmap CM)
\usepackage{textcomp}           % extra text symbols
\usepackage{microtype}          % subtle kerning improvements

% ---- UTF-8 → LaTeX character mappings --------------------------------
% Maps the Unicode characters that occur in TECT body text (legacy from
% the chat-content auto-archival rule, CLAUDE.md §4) to LaTeX-safe
% commands. Without these declarations, T1+inputenc rejects the chars
% with "Unicode character X not set up for use with LaTeX".
% Adding declarations centrally (here) is preferable to per-file edits.
\DeclareUnicodeCharacter{00A7}{\S}              % §  (section sign)
\DeclareUnicodeCharacter{00B2}{\ensuremath{^{2}}}        % ²
\DeclareUnicodeCharacter{00B3}{\ensuremath{^{3}}}        % ³
\DeclareUnicodeCharacter{00B5}{\ensuremath{\mu}}         % µ (micro)
\DeclareUnicodeCharacter{00B7}{\ensuremath{\cdot}}       % · (middle dot)
\DeclareUnicodeCharacter{00D7}{\ensuremath{\times}}      % ×
\DeclareUnicodeCharacter{00E4}{\"a}             % ä
\DeclareUnicodeCharacter{00E9}{\'e}             % é
\DeclareUnicodeCharacter{00F6}{\"o}             % ö
\DeclareUnicodeCharacter{00FC}{\"u}             % ü
\DeclareUnicodeCharacter{010C}{\v{C}}           % Č
\DeclareUnicodeCharacter{010D}{\v{c}}           % č
\DeclareUnicodeCharacter{0394}{\ensuremath{\Delta}}      % Δ
\DeclareUnicodeCharacter{03A3}{\ensuremath{\Sigma}}      % Σ
\DeclareUnicodeCharacter{03A9}{\ensuremath{\Omega}}      % Ω
\DeclareUnicodeCharacter{03B1}{\ensuremath{\alpha}}      % α
\DeclareUnicodeCharacter{03B2}{\ensuremath{\beta}}       % β
\DeclareUnicodeCharacter{03B3}{\ensuremath{\gamma}}      % γ
\DeclareUnicodeCharacter{03B4}{\ensuremath{\delta}}      % δ
\DeclareUnicodeCharacter{03B5}{\ensuremath{\varepsilon}} % ε
\DeclareUnicodeCharacter{03BA}{\ensuremath{\kappa}}      % κ
\DeclareUnicodeCharacter{03BB}{\ensuremath{\lambda}}     % λ
\DeclareUnicodeCharacter{03BC}{\ensuremath{\mu}}         % μ
\DeclareUnicodeCharacter{03BD}{\ensuremath{\nu}}         % ν
\DeclareUnicodeCharacter{03C0}{\ensuremath{\pi}}         % π
\DeclareUnicodeCharacter{03C1}{\ensuremath{\rho}}        % ρ
\DeclareUnicodeCharacter{03C3}{\ensuremath{\sigma}}      % σ
\DeclareUnicodeCharacter{03C4}{\ensuremath{\tau}}        % τ
\DeclareUnicodeCharacter{03C6}{\ensuremath{\varphi}}     % φ
\DeclareUnicodeCharacter{03C7}{\ensuremath{\chi}}        % χ
\DeclareUnicodeCharacter{03C8}{\ensuremath{\psi}}        % ψ
\DeclareUnicodeCharacter{03C9}{\ensuremath{\omega}}      % ω
\DeclareUnicodeCharacter{2013}{--}              % – (en dash)
\DeclareUnicodeCharacter{2014}{---}             % — (em dash)
\DeclareUnicodeCharacter{2018}{`}               % ‘
\DeclareUnicodeCharacter{2019}{'}               % ’
\DeclareUnicodeCharacter{201C}{``}              % “
\DeclareUnicodeCharacter{201D}{''}              % ”
\DeclareUnicodeCharacter{2070}{\ensuremath{^{0}}}        % ⁰
\DeclareUnicodeCharacter{2071}{\ensuremath{^{i}}}        % ⁱ
\DeclareUnicodeCharacter{2122}{\textsuperscript{TM}}     % ™
\DeclareUnicodeCharacter{2126}{\ensuremath{\Omega}}      % Ω (ohm sign)
\DeclareUnicodeCharacter{210F}{\ensuremath{\hbar}}       % ℏ
\DeclareUnicodeCharacter{2115}{\ensuremath{\mathbb{N}}}  % ℕ
\DeclareUnicodeCharacter{2102}{\ensuremath{\mathbb{C}}}  % ℂ
\DeclareUnicodeCharacter{211D}{\ensuremath{\mathbb{R}}}  % ℝ
\DeclareUnicodeCharacter{2124}{\ensuremath{\mathbb{Z}}}  % ℤ
\DeclareUnicodeCharacter{211A}{\ensuremath{\mathbb{Q}}}  % ℚ
\DeclareUnicodeCharacter{2192}{\ensuremath{\to}}         % →
\DeclareUnicodeCharacter{21D2}{\ensuremath{\Rightarrow}} % ⇒
\DeclareUnicodeCharacter{2200}{\ensuremath{\forall}}     % ∀
\DeclareUnicodeCharacter{2203}{\ensuremath{\exists}}     % ∃
\DeclareUnicodeCharacter{2208}{\ensuremath{\in}}         % ∈
\DeclareUnicodeCharacter{2209}{\ensuremath{\notin}}      % ∉
\DeclareUnicodeCharacter{2227}{\ensuremath{\wedge}}      % ∧
\DeclareUnicodeCharacter{2228}{\ensuremath{\vee}}        % ∨
\DeclareUnicodeCharacter{2229}{\ensuremath{\cap}}        % ∩
\DeclareUnicodeCharacter{222A}{\ensuremath{\cup}}        % ∪
\DeclareUnicodeCharacter{2245}{\ensuremath{\cong}}       % ≅
\DeclareUnicodeCharacter{2248}{\ensuremath{\approx}}     % ≈
\DeclareUnicodeCharacter{2260}{\ensuremath{\neq}}        % ≠
\DeclareUnicodeCharacter{2264}{\ensuremath{\leq}}        % ≤
\DeclareUnicodeCharacter{2265}{\ensuremath{\geq}}        % ≥
\DeclareUnicodeCharacter{22C5}{\ensuremath{\cdot}}       % ⋅
\DeclareUnicodeCharacter{2295}{\ensuremath{\oplus}}      % ⊕
\DeclareUnicodeCharacter{2297}{\ensuremath{\otimes}}     % ⊗

% ---- Mathematics & symbols ------------------------------------------
\usepackage{amsmath,amssymb,bm,mathrsfs,mathtools}
\usepackage{amsthm}             % theorem environments (load BEFORE hyperref)

% ---- Theorem environments (shared counter set, TECT canonical) -------
% All theorems / lemmas / propositions / corollaries / remarks share a
% single counter so cross-references read consistently across a paper.
% Definitions get a separate counter (definitions are numbered in their
% own sequence by editorial convention).
\newtheorem{theorem}{Theorem}
\newtheorem{lemma}[theorem]{Lemma}
\newtheorem{proposition}[theorem]{Proposition}
\newtheorem{corollary}[theorem]{Corollary}

\theoremstyle{remark}
\newtheorem{remark}[theorem]{Remark}
\newtheorem{example}[theorem]{Example}

\theoremstyle{definition}
\newtheorem{definition}[theorem]{Definition}
\newtheorem{conjecture}[theorem]{Conjecture}
\newtheorem{hypothesis}[theorem]{Hypothesis}

% ---- Tables / floats / graphics --------------------------------------
\usepackage{booktabs}           % \toprule / \midrule / \bottomrule
\usepackage{array}
\usepackage{graphicx}
\usepackage{enumitem}           % \begin{itemize}[leftmargin=...] options
\usepackage{hyphenat}           % \hyp{} for breakable hyphens in \texttt
\usepackage{seqsplit}           % \seqsplit{} for long unbreakable strings
\sloppy                         % allow stretchier inter-word spacing to
                                % reduce overfull \hbox warnings on long URLs
                                % and long \texttt tokens (paragraph-level).
\emergencystretch=3em           % last-resort hbox stretch budget

% ---- Cross-references (hyperref last among the standard packages,
%      cleveref AFTER hyperref) ----------------------------------------
\usepackage[%
  colorlinks=true,%
  linkcolor=blue,%
  citecolor=blue,%
  urlcolor=blue,%
  breaklinks=true,%
  bookmarks=true,%
  bookmarksnumbered=true%
]{hyperref}
\usepackage[capitalise,nameinlink]{cleveref}

% ---- TECT-canonical author block macros -----------------------------
% (defined here so every paper has the same author / affiliation style)
\newcommand{\tectauthor}{%
  \author{Jusang Lee}%
  \email{jtkor@outlook.com}%
  \affiliation{Independent researcher, Republic of Korea}%
}

% ---- TECT-canonical archive footer ----------------------------------
\newcommand{\tectarchivefooter}{%
  \par\medskip
  \noindent\small\textit{Canonical archive}:
  \url{https://github.com/TECT-OpenPhysics/TECT}.\par%
}

% ---- TECT TODO / AUDIT-FLAG / HONEST-SCOPE inline marks --------------
% Used in drafts to highlight pending audit items without disturbing
% the body text style. Suppress via \tectsuppressmarks (define empty).
\providecommand{\tectauditflag}[1]{\textcolor{red}{\textbf{[AUDIT-FLAG: #1]}}}
\providecommand{\tecttodo}[1]{\textcolor{orange}{\textbf{[TODO: #1]}}}
\providecommand{\tecthonestscope}[1]{\textit{Honest scope: #1.}}

% =====================================================================
% End of TECT canonical preamble
% =====================================================================

%   \begin{document}
%   ...
%   \end{document}
%
% Build (every paper): `pdflatex Paper-NN && pdflatex Paper-NN`
% (the second pass resolves \ref{} / \cref{} forward references).
% A bash/PowerShell wrapper that automates this is at
% Docs/papers/papers/_shared/build-paper.sh / build-paper.ps1.
% =====================================================================

% ---- Document class --------------------------------------------------
% PRD class option of revtex4-2 with [reprint,onecolumn] gives the same
% APS heading / by-line / abstract style as PRL but in a wider single
% column that handles long display equations and theorem environments
% without column-break artefacts. nofootinbib keeps citations out of
% footnotes. The 11pt option restores standard font metrics; with T1
% fontenc + lmodern this gives non-bitmap glyphs.
\documentclass[%
  aps,prd,reprint,onecolumn,nofootinbib,11pt,%
  amsmath,amssymb,floatfix,longbibliography%
]{revtex4-2}

% ---- Encoding & fonts ------------------------------------------------
\usepackage[T1]{fontenc}        % proper 8-bit font encoding (no font corruption)
\usepackage[utf8]{inputenc}     % allow UTF-8 source
\usepackage{lmodern}            % scalable Latin Modern (replaces bitmap CM)
\usepackage{textcomp}           % extra text symbols
\usepackage{microtype}          % subtle kerning improvements

% ---- UTF-8 → LaTeX character mappings --------------------------------
% Maps the Unicode characters that occur in TECT body text (legacy from
% the chat-content auto-archival rule, CLAUDE.md §4) to LaTeX-safe
% commands. Without these declarations, T1+inputenc rejects the chars
% with "Unicode character X not set up for use with LaTeX".
% Adding declarations centrally (here) is preferable to per-file edits.
\DeclareUnicodeCharacter{00A7}{\S}              % §  (section sign)
\DeclareUnicodeCharacter{00B2}{\ensuremath{^{2}}}        % ²
\DeclareUnicodeCharacter{00B3}{\ensuremath{^{3}}}        % ³
\DeclareUnicodeCharacter{00B5}{\ensuremath{\mu}}         % µ (micro)
\DeclareUnicodeCharacter{00B7}{\ensuremath{\cdot}}       % · (middle dot)
\DeclareUnicodeCharacter{00D7}{\ensuremath{\times}}      % ×
\DeclareUnicodeCharacter{00E4}{\"a}             % ä
\DeclareUnicodeCharacter{00E9}{\'e}             % é
\DeclareUnicodeCharacter{00F6}{\"o}             % ö
\DeclareUnicodeCharacter{00FC}{\"u}             % ü
\DeclareUnicodeCharacter{010C}{\v{C}}           % Č
\DeclareUnicodeCharacter{010D}{\v{c}}           % č
\DeclareUnicodeCharacter{0394}{\ensuremath{\Delta}}      % Δ
\DeclareUnicodeCharacter{03A3}{\ensuremath{\Sigma}}      % Σ
\DeclareUnicodeCharacter{03A9}{\ensuremath{\Omega}}      % Ω
\DeclareUnicodeCharacter{03B1}{\ensuremath{\alpha}}      % α
\DeclareUnicodeCharacter{03B2}{\ensuremath{\beta}}       % β
\DeclareUnicodeCharacter{03B3}{\ensuremath{\gamma}}      % γ
\DeclareUnicodeCharacter{03B4}{\ensuremath{\delta}}      % δ
\DeclareUnicodeCharacter{03B5}{\ensuremath{\varepsilon}} % ε
\DeclareUnicodeCharacter{03BA}{\ensuremath{\kappa}}      % κ
\DeclareUnicodeCharacter{03BB}{\ensuremath{\lambda}}     % λ
\DeclareUnicodeCharacter{03BC}{\ensuremath{\mu}}         % μ
\DeclareUnicodeCharacter{03BD}{\ensuremath{\nu}}         % ν
\DeclareUnicodeCharacter{03C0}{\ensuremath{\pi}}         % π
\DeclareUnicodeCharacter{03C1}{\ensuremath{\rho}}        % ρ
\DeclareUnicodeCharacter{03C3}{\ensuremath{\sigma}}      % σ
\DeclareUnicodeCharacter{03C4}{\ensuremath{\tau}}        % τ
\DeclareUnicodeCharacter{03C6}{\ensuremath{\varphi}}     % φ
\DeclareUnicodeCharacter{03C7}{\ensuremath{\chi}}        % χ
\DeclareUnicodeCharacter{03C8}{\ensuremath{\psi}}        % ψ
\DeclareUnicodeCharacter{03C9}{\ensuremath{\omega}}      % ω
\DeclareUnicodeCharacter{2013}{--}              % – (en dash)
\DeclareUnicodeCharacter{2014}{---}             % — (em dash)
\DeclareUnicodeCharacter{2018}{`}               % ‘
\DeclareUnicodeCharacter{2019}{'}               % ’
\DeclareUnicodeCharacter{201C}{``}              % “
\DeclareUnicodeCharacter{201D}{''}              % ”
\DeclareUnicodeCharacter{2070}{\ensuremath{^{0}}}        % ⁰
\DeclareUnicodeCharacter{2071}{\ensuremath{^{i}}}        % ⁱ
\DeclareUnicodeCharacter{2122}{\textsuperscript{TM}}     % ™
\DeclareUnicodeCharacter{2126}{\ensuremath{\Omega}}      % Ω (ohm sign)
\DeclareUnicodeCharacter{210F}{\ensuremath{\hbar}}       % ℏ
\DeclareUnicodeCharacter{2115}{\ensuremath{\mathbb{N}}}  % ℕ
\DeclareUnicodeCharacter{2102}{\ensuremath{\mathbb{C}}}  % ℂ
\DeclareUnicodeCharacter{211D}{\ensuremath{\mathbb{R}}}  % ℝ
\DeclareUnicodeCharacter{2124}{\ensuremath{\mathbb{Z}}}  % ℤ
\DeclareUnicodeCharacter{211A}{\ensuremath{\mathbb{Q}}}  % ℚ
\DeclareUnicodeCharacter{2192}{\ensuremath{\to}}         % →
\DeclareUnicodeCharacter{21D2}{\ensuremath{\Rightarrow}} % ⇒
\DeclareUnicodeCharacter{2200}{\ensuremath{\forall}}     % ∀
\DeclareUnicodeCharacter{2203}{\ensuremath{\exists}}     % ∃
\DeclareUnicodeCharacter{2208}{\ensuremath{\in}}         % ∈
\DeclareUnicodeCharacter{2209}{\ensuremath{\notin}}      % ∉
\DeclareUnicodeCharacter{2227}{\ensuremath{\wedge}}      % ∧
\DeclareUnicodeCharacter{2228}{\ensuremath{\vee}}        % ∨
\DeclareUnicodeCharacter{2229}{\ensuremath{\cap}}        % ∩
\DeclareUnicodeCharacter{222A}{\ensuremath{\cup}}        % ∪
\DeclareUnicodeCharacter{2245}{\ensuremath{\cong}}       % ≅
\DeclareUnicodeCharacter{2248}{\ensuremath{\approx}}     % ≈
\DeclareUnicodeCharacter{2260}{\ensuremath{\neq}}        % ≠
\DeclareUnicodeCharacter{2264}{\ensuremath{\leq}}        % ≤
\DeclareUnicodeCharacter{2265}{\ensuremath{\geq}}        % ≥
\DeclareUnicodeCharacter{22C5}{\ensuremath{\cdot}}       % ⋅
\DeclareUnicodeCharacter{2295}{\ensuremath{\oplus}}      % ⊕
\DeclareUnicodeCharacter{2297}{\ensuremath{\otimes}}     % ⊗

% ---- Mathematics & symbols ------------------------------------------
\usepackage{amsmath,amssymb,bm,mathrsfs,mathtools}
\usepackage{amsthm}             % theorem environments (load BEFORE hyperref)

% ---- Theorem environments (shared counter set, TECT canonical) -------
% All theorems / lemmas / propositions / corollaries / remarks share a
% single counter so cross-references read consistently across a paper.
% Definitions get a separate counter (definitions are numbered in their
% own sequence by editorial convention).
\newtheorem{theorem}{Theorem}
\newtheorem{lemma}[theorem]{Lemma}
\newtheorem{proposition}[theorem]{Proposition}
\newtheorem{corollary}[theorem]{Corollary}

\theoremstyle{remark}
\newtheorem{remark}[theorem]{Remark}
\newtheorem{example}[theorem]{Example}

\theoremstyle{definition}
\newtheorem{definition}[theorem]{Definition}
\newtheorem{conjecture}[theorem]{Conjecture}
\newtheorem{hypothesis}[theorem]{Hypothesis}

% ---- Tables / floats / graphics --------------------------------------
\usepackage{booktabs}           % \toprule / \midrule / \bottomrule
\usepackage{array}
\usepackage{graphicx}
\usepackage{enumitem}           % \begin{itemize}[leftmargin=...] options
\usepackage{hyphenat}           % \hyp{} for breakable hyphens in \texttt
\usepackage{seqsplit}           % \seqsplit{} for long unbreakable strings
\sloppy                         % allow stretchier inter-word spacing to
                                % reduce overfull \hbox warnings on long URLs
                                % and long \texttt tokens (paragraph-level).
\emergencystretch=3em           % last-resort hbox stretch budget

% ---- Cross-references (hyperref last among the standard packages,
%      cleveref AFTER hyperref) ----------------------------------------
\usepackage[%
  colorlinks=true,%
  linkcolor=blue,%
  citecolor=blue,%
  urlcolor=blue,%
  breaklinks=true,%
  bookmarks=true,%
  bookmarksnumbered=true%
]{hyperref}
\usepackage[capitalise,nameinlink]{cleveref}

% ---- TECT-canonical author block macros -----------------------------
% (defined here so every paper has the same author / affiliation style)
\newcommand{\tectauthor}{%
  \author{Jusang Lee}%
  \email{jtkor@outlook.com}%
  \affiliation{Independent researcher, Republic of Korea}%
}

% ---- TECT-canonical archive footer ----------------------------------
\newcommand{\tectarchivefooter}{%
  \par\medskip
  \noindent\small\textit{Canonical archive}:
  \url{https://github.com/TECT-OpenPhysics/TECT}.\par%
}

% ---- TECT TODO / AUDIT-FLAG / HONEST-SCOPE inline marks --------------
% Used in drafts to highlight pending audit items without disturbing
% the body text style. Suppress via \tectsuppressmarks (define empty).
\providecommand{\tectauditflag}[1]{\textcolor{red}{\textbf{[AUDIT-FLAG: #1]}}}
\providecommand{\tecttodo}[1]{\textcolor{orange}{\textbf{[TODO: #1]}}}
\providecommand{\tecthonestscope}[1]{\textit{Honest scope: #1.}}

% =====================================================================
% End of TECT canonical preamble
% =====================================================================

%   \begin{document}
%   ...
%   \end{document}
%
% Build (every paper): `pdflatex Paper-NN && pdflatex Paper-NN`
% (the second pass resolves \ref{} / \cref{} forward references).
% A bash/PowerShell wrapper that automates this is at
% Docs/papers/papers/_shared/build-paper.sh / build-paper.ps1.
% =====================================================================

% ---- Document class --------------------------------------------------
% PRD class option of revtex4-2 with [reprint,onecolumn] gives the same
% APS heading / by-line / abstract style as PRL but in a wider single
% column that handles long display equations and theorem environments
% without column-break artefacts. nofootinbib keeps citations out of
% footnotes. The 11pt option restores standard font metrics; with T1
% fontenc + lmodern this gives non-bitmap glyphs.
\documentclass[%
  aps,prd,reprint,onecolumn,nofootinbib,11pt,%
  amsmath,amssymb,floatfix,longbibliography%
]{revtex4-2}

% ---- Encoding & fonts ------------------------------------------------
\usepackage[T1]{fontenc}        % proper 8-bit font encoding (no font corruption)
\usepackage[utf8]{inputenc}     % allow UTF-8 source
\usepackage{lmodern}            % scalable Latin Modern (replaces bitmap CM)
\usepackage{textcomp}           % extra text symbols
\usepackage{microtype}          % subtle kerning improvements

% ---- UTF-8 → LaTeX character mappings --------------------------------
% Maps the Unicode characters that occur in TECT body text (legacy from
% the chat-content auto-archival rule, CLAUDE.md §4) to LaTeX-safe
% commands. Without these declarations, T1+inputenc rejects the chars
% with "Unicode character X not set up for use with LaTeX".
% Adding declarations centrally (here) is preferable to per-file edits.
\DeclareUnicodeCharacter{00A7}{\S}              % §  (section sign)
\DeclareUnicodeCharacter{00B2}{\ensuremath{^{2}}}        % ²
\DeclareUnicodeCharacter{00B3}{\ensuremath{^{3}}}        % ³
\DeclareUnicodeCharacter{00B5}{\ensuremath{\mu}}         % µ (micro)
\DeclareUnicodeCharacter{00B7}{\ensuremath{\cdot}}       % · (middle dot)
\DeclareUnicodeCharacter{00D7}{\ensuremath{\times}}      % ×
\DeclareUnicodeCharacter{00E4}{\"a}             % ä
\DeclareUnicodeCharacter{00E9}{\'e}             % é
\DeclareUnicodeCharacter{00F6}{\"o}             % ö
\DeclareUnicodeCharacter{00FC}{\"u}             % ü
\DeclareUnicodeCharacter{010C}{\v{C}}           % Č
\DeclareUnicodeCharacter{010D}{\v{c}}           % č
\DeclareUnicodeCharacter{0394}{\ensuremath{\Delta}}      % Δ
\DeclareUnicodeCharacter{03A3}{\ensuremath{\Sigma}}      % Σ
\DeclareUnicodeCharacter{03A9}{\ensuremath{\Omega}}      % Ω
\DeclareUnicodeCharacter{03B1}{\ensuremath{\alpha}}      % α
\DeclareUnicodeCharacter{03B2}{\ensuremath{\beta}}       % β
\DeclareUnicodeCharacter{03B3}{\ensuremath{\gamma}}      % γ
\DeclareUnicodeCharacter{03B4}{\ensuremath{\delta}}      % δ
\DeclareUnicodeCharacter{03B5}{\ensuremath{\varepsilon}} % ε
\DeclareUnicodeCharacter{03BA}{\ensuremath{\kappa}}      % κ
\DeclareUnicodeCharacter{03BB}{\ensuremath{\lambda}}     % λ
\DeclareUnicodeCharacter{03BC}{\ensuremath{\mu}}         % μ
\DeclareUnicodeCharacter{03BD}{\ensuremath{\nu}}         % ν
\DeclareUnicodeCharacter{03C0}{\ensuremath{\pi}}         % π
\DeclareUnicodeCharacter{03C1}{\ensuremath{\rho}}        % ρ
\DeclareUnicodeCharacter{03C3}{\ensuremath{\sigma}}      % σ
\DeclareUnicodeCharacter{03C4}{\ensuremath{\tau}}        % τ
\DeclareUnicodeCharacter{03C6}{\ensuremath{\varphi}}     % φ
\DeclareUnicodeCharacter{03C7}{\ensuremath{\chi}}        % χ
\DeclareUnicodeCharacter{03C8}{\ensuremath{\psi}}        % ψ
\DeclareUnicodeCharacter{03C9}{\ensuremath{\omega}}      % ω
\DeclareUnicodeCharacter{2013}{--}              % – (en dash)
\DeclareUnicodeCharacter{2014}{---}             % — (em dash)
\DeclareUnicodeCharacter{2018}{`}               % ‘
\DeclareUnicodeCharacter{2019}{'}               % ’
\DeclareUnicodeCharacter{201C}{``}              % “
\DeclareUnicodeCharacter{201D}{''}              % ”
\DeclareUnicodeCharacter{2070}{\ensuremath{^{0}}}        % ⁰
\DeclareUnicodeCharacter{2071}{\ensuremath{^{i}}}        % ⁱ
\DeclareUnicodeCharacter{2122}{\textsuperscript{TM}}     % ™
\DeclareUnicodeCharacter{2126}{\ensuremath{\Omega}}      % Ω (ohm sign)
\DeclareUnicodeCharacter{210F}{\ensuremath{\hbar}}       % ℏ
\DeclareUnicodeCharacter{2115}{\ensuremath{\mathbb{N}}}  % ℕ
\DeclareUnicodeCharacter{2102}{\ensuremath{\mathbb{C}}}  % ℂ
\DeclareUnicodeCharacter{211D}{\ensuremath{\mathbb{R}}}  % ℝ
\DeclareUnicodeCharacter{2124}{\ensuremath{\mathbb{Z}}}  % ℤ
\DeclareUnicodeCharacter{211A}{\ensuremath{\mathbb{Q}}}  % ℚ
\DeclareUnicodeCharacter{2192}{\ensuremath{\to}}         % →
\DeclareUnicodeCharacter{21D2}{\ensuremath{\Rightarrow}} % ⇒
\DeclareUnicodeCharacter{2200}{\ensuremath{\forall}}     % ∀
\DeclareUnicodeCharacter{2203}{\ensuremath{\exists}}     % ∃
\DeclareUnicodeCharacter{2208}{\ensuremath{\in}}         % ∈
\DeclareUnicodeCharacter{2209}{\ensuremath{\notin}}      % ∉
\DeclareUnicodeCharacter{2227}{\ensuremath{\wedge}}      % ∧
\DeclareUnicodeCharacter{2228}{\ensuremath{\vee}}        % ∨
\DeclareUnicodeCharacter{2229}{\ensuremath{\cap}}        % ∩
\DeclareUnicodeCharacter{222A}{\ensuremath{\cup}}        % ∪
\DeclareUnicodeCharacter{2245}{\ensuremath{\cong}}       % ≅
\DeclareUnicodeCharacter{2248}{\ensuremath{\approx}}     % ≈
\DeclareUnicodeCharacter{2260}{\ensuremath{\neq}}        % ≠
\DeclareUnicodeCharacter{2264}{\ensuremath{\leq}}        % ≤
\DeclareUnicodeCharacter{2265}{\ensuremath{\geq}}        % ≥
\DeclareUnicodeCharacter{22C5}{\ensuremath{\cdot}}       % ⋅
\DeclareUnicodeCharacter{2295}{\ensuremath{\oplus}}      % ⊕
\DeclareUnicodeCharacter{2297}{\ensuremath{\otimes}}     % ⊗

% ---- Mathematics & symbols ------------------------------------------
\usepackage{amsmath,amssymb,bm,mathrsfs,mathtools}
\usepackage{amsthm}             % theorem environments (load BEFORE hyperref)

% ---- Theorem environments (shared counter set, TECT canonical) -------
% All theorems / lemmas / propositions / corollaries / remarks share a
% single counter so cross-references read consistently across a paper.
% Definitions get a separate counter (definitions are numbered in their
% own sequence by editorial convention).
\newtheorem{theorem}{Theorem}
\newtheorem{lemma}[theorem]{Lemma}
\newtheorem{proposition}[theorem]{Proposition}
\newtheorem{corollary}[theorem]{Corollary}

\theoremstyle{remark}
\newtheorem{remark}[theorem]{Remark}
\newtheorem{example}[theorem]{Example}

\theoremstyle{definition}
\newtheorem{definition}[theorem]{Definition}
\newtheorem{conjecture}[theorem]{Conjecture}
\newtheorem{hypothesis}[theorem]{Hypothesis}

% ---- Tables / floats / graphics --------------------------------------
\usepackage{booktabs}           % \toprule / \midrule / \bottomrule
\usepackage{array}
\usepackage{graphicx}
\usepackage{enumitem}           % \begin{itemize}[leftmargin=...] options
\usepackage{hyphenat}           % \hyp{} for breakable hyphens in \texttt
\usepackage{seqsplit}           % \seqsplit{} for long unbreakable strings
\sloppy                         % allow stretchier inter-word spacing to
                                % reduce overfull \hbox warnings on long URLs
                                % and long \texttt tokens (paragraph-level).
\emergencystretch=3em           % last-resort hbox stretch budget

% ---- Cross-references (hyperref last among the standard packages,
%      cleveref AFTER hyperref) ----------------------------------------
\usepackage[%
  colorlinks=true,%
  linkcolor=blue,%
  citecolor=blue,%
  urlcolor=blue,%
  breaklinks=true,%
  bookmarks=true,%
  bookmarksnumbered=true%
]{hyperref}
\usepackage[capitalise,nameinlink]{cleveref}

% ---- TECT-canonical author block macros -----------------------------
% (defined here so every paper has the same author / affiliation style)
\newcommand{\tectauthor}{%
  \author{Jusang Lee}%
  \email{jtkor@outlook.com}%
  \affiliation{Independent researcher, Republic of Korea}%
}

% ---- TECT-canonical archive footer ----------------------------------
\newcommand{\tectarchivefooter}{%
  \par\medskip
  \noindent\small\textit{Canonical archive}:
  \url{https://github.com/TECT-OpenPhysics/TECT}.\par%
}

% ---- TECT TODO / AUDIT-FLAG / HONEST-SCOPE inline marks --------------
% Used in drafts to highlight pending audit items without disturbing
% the body text style. Suppress via \tectsuppressmarks (define empty).
\providecommand{\tectauditflag}[1]{\textcolor{red}{\textbf{[AUDIT-FLAG: #1]}}}
\providecommand{\tecttodo}[1]{\textcolor{orange}{\textbf{[TODO: #1]}}}
\providecommand{\tecthonestscope}[1]{\textit{Honest scope: #1.}}

% =====================================================================
% End of TECT canonical preamble
% =====================================================================


\begin{document}

\title{Pillar~4 of the Topological Energy Condensate Theory:
$\mathrm{SU}(5)\times\mathrm{U}(1)$ gauge emergence from BCC topological
defects on $\Sigma_0 = \mathbb{P}^1\times\mathbb{P}^1$, Cartan-forcing
algebraic closure, triple verification, and the F-GAP4 verdict shell}

\author{Jusang Lee}
\email{jtkor@outlook.com}
\affiliation{Independent researcher, Republic of Korea}

\date{\today}

\begin{abstract}
We present the Pillar-4 atomic theorem of the Topological Energy
Condensate Theory (TECT): the emergence of an
$\mathrm{SU}(5) \times \mathrm{U}(1)$ gauge structure from the
topological-defect sector of the BCC condensate, formulated on the
canonical Kähler base $\Sigma_0 = \mathbb{P}^1 \times \mathbb{P}^1$.
The result is the conjunction of three sub-task theorems: (i)
existence and properties of an SO(10)-structured BCC-defect bundle
$E \to \Sigma_0$ with explicit four-chart \v{C}ech atlas and lifted
transition functions; (ii) representation-theoretic decomposition of
$E$ under the cubic point group $O_h$ with emergent $\mathrm{SU}(5)$
breaking; (iii) Cartan-forcing algebraic closure of the
$\mathrm{SO}(10) \to \mathrm{SU}(5) \times \mathrm{U}(1) \to
\mathrm{SM}$ chain via a first-principles re-derivation of the
$O_h$-character table by Weyl formula and Frobenius reciprocity.
The atomic claim carries the canonical tier
\textbf{T6 PROVED CONDITIONAL} on a thirteen-hypothesis set
$\mathcal{H}_{\rm atomic} = \{\mathrm{H1.1, H2.1\text{--}H2.7, H3.1,
H_{task}}\}$, of which eight are at T6, two at T7 (textbook K\"ahler
and Cartan-forcing), and one (the $H_{\rm task}$ realisation gate
for F-GAP4-DEFECT-MASS, deadline 2026-05-14) is a programme-completion
condition that does NOT downgrade the proof-rigor tier per the
Math302 carve-out. The result has been verified at three independent
levels: (V1) original promotion (Math268--269); (V2) external defense
programme of five hostile-referee attacks (Math270--275), all
DISCHARGED or MITIGATED in the Math276 cross-turn audit (OUTCOME A);
(V3) two third-order audits of the load-bearing Lemmas A and B
(Math277, Math278), both OUTCOME A. The F-GAP4-DEFECT-MASS verdict
shell (Math300 + Math311) maps the defect-mass eigenvalue
$\hat\lambda_{\min}(\Box_E)$ deterministically to a window-verdict
$\mu_{\rm defect} \in [10^{13},\,10^{17}]\,\mathrm{GeV}$ with
solver/grid certificates: PASS upgrades the canonical realisation
tier, FAIL activates the Math246 contingency routes (B = twisted
compactification, C = orbifold cover, D = heterotic Calabi--Yau
3-fold), DEFER triggers Math246 Task Q escalation. The verdict shell
is explicitly a window check, not a $\pm 5\%$ precision prediction.
The earlier $\mathbb{CP}^2$-based version of the programme is
preserved in the companion top-impact paper TI-1 as a background
algebro-geometric correction to the spin-c index calculation; the
present Pillar-4 paper operates on the canonical $\Sigma_0$ base
introduced in Math288 and explicitly constructed in Math305.
\end{abstract}

\maketitle

% =====================================================================
\section{Introduction and structural context}\label{sec:intro}
% =====================================================================

Pillar~4 of the TECT framework is the structural-emergence theorem
that the $\mathrm{SU}(5) \times \mathrm{U}(1)$ gauge sector of the
candidate Standard-Model embedding arises from the topological-defect
sector of the BCC condensate, rather than being imposed as an
external postulate. The present paper reports the canonical Pillar-4
atomic theorem in its current form (Rev 3, 2026-05-03) following the
post-defense final consolidation (Math279) and the explicit
$\Sigma_0$-realisation construction (Math305 + Math288). The result
is conditional on an explicitly enumerated hypothesis set, not
unconditional; the conditional structure is tracked separately at the
proof-rigor level and at the realisation-gate level per the Math302
carve-out.

The structural development is organised in three layers. \emph{First},
the Pillar-4 atomic statement (\S\ref{sec:atomic}) is a conjunction
of three sub-task theorems and a cumulative hypothesis set
$\mathcal{H}_{\rm atomic}$ of thirteen hypotheses; the composite tier
under the internal audit-policy minimum rule is T6 PROVED CONDITIONAL.
\emph{Second}, the explicit $\Sigma_0 = \mathbb{P}^1 \times
\mathbb{P}^1$ \v{C}ech atlas with lifted SO(10) transition functions
is constructed in \S\ref{sec:sigma0} (after Math305). \emph{Third},
the F-GAP4-DEFECT-MASS verdict shell (Math300 + Math311) and the
Math288 defect-mass formula derivation are summarised in
\S\ref{sec:fgap4}: the verdict is a window-shell decision tree, not
a precision prediction.

The historical $\mathbb{CP}^2$-based version of the Pillar-4
programme (Math162 \v{C}ech construction; Math174 explicit $c_2 =
-40$ Chern-class calculation; Math191--192 $c_1, c_2$ from principal
connections) is preserved in the companion top-impact paper TI-1 as
a background correction to the spin-c index formula on
$\mathbb{CP}^2$. The transfer of the topological certificate from
$\mathbb{CP}^2$ to $\Sigma_0$ is established by the Math270 Theorem
270.1 (T7-rigorous splitting-principle certificate transfer).

% =====================================================================
\section{Pillar-4 atomic theorem and the hypothesis
set}\label{sec:atomic}
% =====================================================================

\subsection{The Brazovskii operating point and the
$\Sigma_0$ base}

The TECT condensate carries the Brazovskii free energy
\begin{equation}
\mathcal{F}[\Psi] = \int d^3 x \left[
\tfrac{1}{2}|\nabla \Psi|^2 + \mu^2 |\Psi|^2
+ \tfrac{\lambda}{4}|\Psi|^4
+ \text{higher-order}
\right],
\label{eq:brazovskii}
\end{equation}
at locked parameters $(\mu^2, \lambda, \gamma) = (0.26, -0.43, 1.62)$
in dimensionless lattice units $a_{\rm BCC} = 1$. The Pillar-4
realisation programme operates on the canonical Kähler base
$\Sigma_0 = \mathbb{P}^1 \times \mathbb{P}^1$ introduced in Math288
and constructed explicitly in Math305 (\S\ref{sec:sigma0} below).

\subsection{The three sub-task theorems}

\begin{theorem}[Pillar-4 atomic, conjunction of three sub-tasks]
\label{thm:atomic}
The Pillar-4 atomic statement is the conjunction of:
\begin{enumerate}
\item \textbf{Sub-task 1} (Math162 + Math164 + Math170 + Math305):
existence and properties of a holomorphic vector bundle $E \to
\Sigma_0$ with structure group $\mathrm{SO}(10)/\mathrm{SU}(5)$,
first Chern class $c_1(E) = 0$ in the SO(10) traceless adjoint and
$c_1^{(j)} = 1$ on each $\mathbb{P}^1$ factor ($j = 1,2$), and
moduli space $\pi_1(M_{\rm BCC}) = \{e\}$ (trivial fundamental group
on physically relevant BCC configurations);
\item \textbf{Sub-task 2} (Math238 + Math239 + Math248 + Math266 +
Math267): decomposition of $E$ into irreducible representations
under the cubic point group $O_h$, with an emergent
$\mathrm{SU}(5)$ structure characterising the breaking pattern of
SO(10) under the Cartan-forcing condition $\{\hat{D}_i, \hat{D}_j\}
\sim c_{ij} \mathcal{O}_h$ on the defect space;
\item \textbf{Sub-task 3} (Math229 + Math275, with Math277+Math278
third-order audits): algebraic closure of the symmetry-breaking chain
$\mathrm{SO}(10) \to \mathrm{SU}(5) \times \mathrm{U}(1) \to
\mathrm{SM}$ via representation-theoretic decomposition under
$O_h$, independent of the base geometry, with dual verification
routes (Route~A group-theoretic, Route~B geometric).
\end{enumerate}
The composite tier under the internal audit-policy minimum rule is
\[
T_{\rm Pillar 4} = \min(T_1, T_2, T_3) = \min(T6, T6, T6) = T6
\quad \text{PROVED CONDITIONAL}
\]
on the hypothesis set $\mathcal{H}_{\rm atomic}^{\rm final}$ enumerated
in \S\ref{sec:hypothesis-set} below.
\end{theorem}

\subsection{The thirteen-hypothesis set
$\mathcal{H}_{\rm atomic}^{\rm final}$}\label{sec:hypothesis-set}

The thirteen hypotheses underpinning Theorem~\ref{thm:atomic} are
enumerated below with their individual tiers (per Math279~\S4):

\paragraph{Sub-task 1.}
\begin{itemize}
\item \textbf{H1.1} (BCC fibre-bundle on $\Sigma_0$ via topological
certificate transfer): T6 PROVED CONDITIONAL per Math162 + Math270
Theorem 270.1.
\end{itemize}

\paragraph{Sub-task 2.}
\begin{itemize}
\item \textbf{H2.1} ($\Sigma_0$ is a Kähler base): T7 PROVED
(textbook).
\item \textbf{H2.2} (Math250 separation ansatz $\Psi(x,\mathbf{k}) =
\phi_0(x) \times f_k$): T6 PROVED CONDITIONAL.
\item \textbf{H2.3} (Donaldson--Uhlenbeck--Yau polystability of $E$):
T6 PROVED CONDITIONAL per Math253.
\item \textbf{H2.4} ($O_h$ embedding of the defect spinor
representation, first-principles re-derivation): T6 PROVED
CONDITIONAL per Math263 + Math275.
\item \textbf{H2.5} (BRST Faddeev--Popov determinant lower bound, on
the separable ansatz): T6 PROVED CONDITIONAL per Math260 +
Math274 tier-split (separable T6, full T4).
\item \textbf{H2.6} ($\hbar$ matching functional, structural form):
T6 PROVED CONDITIONAL per Math261 + Math273 STRUCTURAL/PHYSICAL
distinction.
\item \textbf{H2.7} (Higgs-potential stability with electroweak
minimum): T6 PROVED CONDITIONAL per Math257.
\end{itemize}

\paragraph{Sub-task 3.}
\begin{itemize}
\item \textbf{H3.1} (Cartan forcing $\{\hat{D}_i,\hat{D}_j\} =
c_{ij}(x) \mathcal{O}_h$ implies $\mathrm{SU}(5) \times
\mathrm{U}(1)$ up to discrete quotients): T7 PROVED per Math229,
audit-confirmed by the third-order audits Math277 + Math278.
\end{itemize}

\paragraph{Realisation gate (programme-completion condition).}
\begin{itemize}
\item \textbf{H$_{\rm task}$} (Task~\#156 numerical closure of the
$\Sigma_0$ HYM realisation by the 2026-05-14 hard deadline): T2
CONJECTURE (binding pre-registered assumption per Math270 Theorem
270.2 contingency clause).
\end{itemize}

\subsection{Composite-tier rule and the Math302 carve-out}

A literal application of the composite-tier rule
$T_{\rm composite} = \min$ over the thirteen hypothesis tiers gives
\[
\min(\mathrm{T6}, \mathrm{T7}, \mathrm{T6}, \mathrm{T6},
\mathrm{T6}, \mathrm{T6}, \mathrm{T6}, \mathrm{T6}, \mathrm{T7},
\mathrm{T2}) = \mathrm{T2},
\]
because of the H$_{\rm task}$ realisation gate. This literal demotion
is mathematically correct but pedagogically misleading.

The Math302 carve-out (\textit{proof-rigor / realisation-gate
distinction}) resolves this by separating the proof-rigor dimension
(at which the mathematical content H1.1--H3.1 sits at T6/T7) from
the realisation-gate dimension (H$_{\rm task}$, which is a
programme-completion condition, not a defect of mathematical proof).
Under the carve-out:
\begin{equation}
\boxed{\text{Pillar~4 atomic: } T6 \text{ PROVED CONDITIONAL on }
\mathcal{H}_{\rm atomic}^{\rm final}.}
\end{equation}
The ``CONDITIONAL'' label carries the obligation to enumerate all
conditions (\S\ref{sec:hypothesis-set} above); the ``T6'' label
reflects the mathematical content tier; H$_{\rm task}$ is recorded as
a separately tracked realisation risk per the canonical
\texttt{TOE-FACT-SHEET.md} bookkeeping.

% =====================================================================
\section{Explicit $\Sigma_0$ \v{C}ech atlas and lifted
transition functions (Math305)}\label{sec:sigma0}
% =====================================================================

The canonical realisation base
$\Sigma_0 = \mathbb{P}^1 \times \mathbb{P}^1$ has the standard
4-chart product atlas $\{U_{ij}\}_{i,j \in \{0,1\}}$ where each
factor $\mathbb{P}^1$ carries its standard 2-chart cover
$U_0 = \{[z_0 : z_1] : z_0 \neq 0\}$, $U_1 = \{[z_0 : z_1] : z_1 \neq
0\}$ with affine coordinate transition $v = 1/u$ on
$U_0 \cap U_1 = \mathbb{C}^\times$.

\begin{proposition}[Lifted SO(10) transition functions, after
Math305]
\label{prop:transitions}
The SO(10) gauge bundle on $\Sigma_0$ inherits the Math162
$\mathbb{CP}^2$-based transition functions via the product
structure. The four-chart cover supports the six pairwise
transitions
\begin{align}
g_{(0,0),(1,0)}(u, v) &= \exp(i \arg(u)\, T_3),\\
g_{(0,0),(0,1)}(u, v) &= \exp(i \arg(v)\, T_3'),\\
g_{(0,0),(1,1)}(u, v) &= \exp\big(i \arg(u)\, T_3 + i \arg(v)\, T_3'\big),\\
g_{(1,0),(0,1)}(u, v) &= g_{(1,0),(0,0)} \cdot g_{(0,0),(0,1)},\\
g_{(1,0),(1,1)}(u, v) &= g_{(0,0),(0,1)}\big|_{u \to 1/u},\\
g_{(0,1),(1,1)}(u, v) &= g_{(0,0),(1,0)}\big|_{v \to 1/v},
\end{align}
where $T_3, T_3'$ are commuting $\mathrm{U}(1)$ generators in the
SO(10) Cartan subalgebra and $\arg$ is the principal argument on
$(-\pi, \pi]$. The cocycle condition
$g_{ab} \cdot g_{bc} = g_{ac}$ is verified by direct substitution;
the transitivity rows are the resulting consistency constraints.
\end{proposition}

\begin{proposition}[Chern data on $\Sigma_0$, after Math305]
\label{prop:chern}
For the rank-45 SO(10) adjoint bundle on
$\mathbb{P}^1 \times \mathbb{P}^1$ with the transition functions of
Proposition~\ref{prop:transitions}: the first Chern class of each
$\mathrm{U}(1)$ factor is $c_1^{(j)} = 1$ in $H^2(\Sigma_0,
\mathbb{Z})$ ($j = 1, 2$); the total $c_1(E) = \mathrm{tr}(F)
= 0$ in the SO(10) traceless adjoint; the second Chern class
$c_2(E) = -\mathrm{tr}(F \wedge F)/(8\pi^2)$ is convention-dependent
at the generator-trace normalisation level, with order-of-magnitude
estimate $|c_2(E)|\sim 50\text{--}100$ (vs Math174 $c_2 = -40$ on
$\mathbb{CP}^2$). The exact value awaits the explicit
Atiyah--Singer index computation
$Q$-2026-05-29-Math305-c2-Sigma0-Numerical follow-up; this
convention-dependent uncertainty does not affect the atlas
construction itself.
\end{proposition}

% =====================================================================
\section{F-GAP4-DEFECT-MASS verdict shell (Math300 +
Math311)}\label{sec:fgap4}
% =====================================================================

\subsection{Defect-mass formula derivation (after Math288)}

The lowest defect mass on the BCC condensate is identified with the
zero-mode eigenvalue of the covariant bundle Laplacian
$\Box_E = D_A^* D_A$ acting on bundle-valued zero-form sections of
$E$. For a holomorphic bundle with an HYM connection this operator is
non-negative, and the smallest positive eigenvalue encodes the defect
mass:
\begin{equation}
\mu_{\rm defect}^2 \equiv \lambda_{\min}(\Box_E)
= \min \mathrm{spec}(D_A^* D_A).
\end{equation}
The dimensional matching uses $a_{\rm BCC}^{-1} \sim 10^{16}$~GeV
(Math110-AddI, lattice action), which sets the natural scale of the
$\Sigma_0$ realisation programme.

\subsection{Bogomolov-type eigenvalue bound (after Math288 \S4.3)}

For a holomorphic vector bundle $E$ on a Kähler surface
$(X, \omega)$ with $c_1(E) = 0$, a Bogomolov-type inequality bounds
the bundle Laplacian spectrum from below by a quantity involving
$c_2(E)$, the volume of $X$, and the rank:
\begin{equation}
\lambda_{\min}(\Box_E) \;\geq\; C_{\rm geo}
\cdot \frac{c_2(E)}{\mathrm{vol}(\Sigma_0)}
\cdot \min_{x \in \Sigma_0} |\det F_{A^*}(x)|,
\end{equation}
where $C_{\rm geo}$ is a universal geometric constant ($O(1)$ in
natural units). The rigorous derivation of this inequality is the
\texttt{Q-2026-05-01-Math288-Proposition-4-3-1-Rigorous-Proof}
follow-up in the canonical \texttt{OPEN-QUESTIONS.md}.

\subsection{The verdict shell (after Math311)}

The F-GAP4-DEFECT-MASS verdict shell is a deterministic mapping from
the numerical eigenvalue $\hat\lambda_{\min}(\Box_E)$ produced by
Task~\#156 to a window verdict on $\mu_{\rm defect}$:
\begin{equation}
\mu_{\rm defect} \in \big[10^{13},\ 10^{17}\big]\,\mathrm{GeV}
\quad \Longleftrightarrow \quad \text{F-GAP4 PASS},
\label{eq:window}
\end{equation}
subject to additional HYM-solver convergence and grid-resolution
gates per Math300 Definition. The decision tree on the verdict
arriving at the 2026-05-14 hard gate is:
\begin{enumerate}
\item \textbf{PASS}: Math270 Theorem 270.2 upgrades T6 $\to$ T7 on
the certificate transfer; Pillar~4 atomic promotes from
T6 PROVED CONDITIONAL to a stronger conditional tier per the
Math311 \S\ref{sec:outcomes} canonical-record update.
\item \textbf{FAIL}: Math270 Theorem 270.2 is invalidated on
$\Sigma_0$; activate Math246 contingency routes B (twisted
compactification), C (orbifold cover), or D (heterotic Calabi--Yau
3-fold) by the 2026-05-21 hard escalation gate.
\item \textbf{DEFER}: Pillar~4 atomic remains T6 PROVED CONDITIONAL
on the existing $\mathcal{H}_{\rm atomic}^{\rm final}$; Math246
Task Q escalation activates per Math286.
\end{enumerate}

\begin{remark}[The verdict is a window-shell, not a $\pm 5\%$
precision prediction]
\label{rem:window-vs-precision}
The window in Eq.~\eqref{eq:window} is a four-orders-of-magnitude
band $10^{13}\text{--}10^{17}\,\mathrm{GeV}$ together with explicit
solver and grid certificates. It is NOT a $\pm 5\%$ precision
prediction of a single mass value. Earlier wording of the form
``$M_{\Sigma_0} = 1.26 \times 10^{16}\,\mathrm{GeV}$ within
$\pm 5\%$'' is incompatible with the canonical Math311 verdict shell
and was downgraded in the present revision per the Math314-AddD
audit and CLAUDE.md \S15.6 rule \#7 (window-verdict vs
precision-prediction distinction).
\end{remark}

% =====================================================================
\section{Triple verification of the Pillar-4 atomic
theorem}\label{sec:verification}
% =====================================================================

\subsection{V1 — Original promotion (Math268 + Math269)}

The original promotion of Pillar~4 atomic from T3 PROOF SKETCH
(Math265, Turn 36) to T6 PROVED CONDITIONAL (Math268 Turn 39,
Math269 Turn 40 consolidation) applied the composite-tier rule
$T_{\rm Pillar4} = \min(T_1, T_2, T_3) = T6$ to the three sub-task
theorems individually established at T6.

\subsection{V2 — External defense programme (Math270--275 +
Math276 cross-turn audit)}

Five hostile-referee attacks were launched against the V1 promotion
(Math270--275); each was DISCHARGED or MITIGATED in the Math276
cross-turn audit (Turn 47), which delivered OUTCOME~A (Pillar~4
atomic T6 PROVED CONDITIONAL externally robust). The five attacks
and their resolutions:
\begin{itemize}
\item \textbf{Attack \#1} (Math270--271, cross-base coherence,
$\Sigma_0$ vs $\mathbb{CP}^2$): MITIGATED via Math270
Theorem~270.1 (T7-rigorous topological certificate transfer) +
Theorem~270.2 (T6 conditional explicit $\Sigma_0$ realisation) +
Task~\#156 contingency binding.
\item \textbf{Attack \#2} (Math272, scope reinterpretation,
algebraic vs geometric vs physical): DISCHARGED via documentary
evidence that the original Pillar-4 scope is algebraic (component
A: breaking-chain pattern), with seven cited canonical sources.
\item \textbf{Attack \#3} (Math273, H2.6 semantic instability,
$\hbar$ matching is structural not physical): DISCHARGED via
explicit STRUCTURAL/PHYSICAL layer distinction.
\item \textbf{Attack \#4} (Math274, H2.5 separable-branch only):
DISCHARGED via honest tier-split (H2.5-separable T6, H2.5-full T4;
Pillar~4 atomic uses separable branch only); Task~\#157 opened for
the H2.5-full T6 upgrade.
\item \textbf{Attack \#5} (Math275, H2.4 literature dependence,
character-table not first-principles): DISCHARGED via
first-principles Weyl-formula + Frobenius-reciprocity re-derivation
of the $O_h$ irreducible decomposition.
\end{itemize}

\subsection{V3 — Third-order audits of Lemmas A and B (Math277 +
Math278)}

The two load-bearing foundational lemmas were independently
audited at the third-order level: Math277 (Turn 48, Lemma~A =
Math220-AddB, the BRST separable-branch determinant lower bound) and
Math278 (Turn 49, Lemma~B = Math221-AddC, the
Donaldson--Uhlenbeck--Yau-applicability technical condition). Both
audits delivered OUTCOME~A: the underlying T6 PROVED CONDITIONAL
claims are mathematically rigorous and audit-confirmed.

\subsection{Combined verification status}

Pillar~4 atomic is therefore triply verified: original proof (V1),
externally robust against five hostile attacks (V2), with
foundational lemmas independently audit-confirmed at the third-order
level (V3). The cumulative verification record is recorded in the
canonical \texttt{Docs/status/TOE-FACT-SHEET.md} as the Stage-1
scorecard entry for Pillar~4.

% =====================================================================
\section{Discussion: scope, what is and is not closed,
contingencies}\label{sec:discussion}
% =====================================================================

\paragraph{What the present paper establishes.}
The Pillar-4 atomic theorem (Theorem~\ref{thm:atomic}) is the
T6-PROVED-CONDITIONAL claim that an
$\mathrm{SU}(5) \times \mathrm{U}(1)$ gauge structure emerges from
the BCC topological-defect sector on the canonical Kähler base
$\Sigma_0 = \mathbb{P}^1 \times \mathbb{P}^1$, conditional on the
explicitly enumerated thirteen-hypothesis set
$\mathcal{H}_{\rm atomic}^{\rm final}$ of
\S\ref{sec:hypothesis-set}. The result is supported by the explicit
$\Sigma_0$ \v{C}ech atlas (\S\ref{sec:sigma0}), the
F-GAP4 verdict shell (\S\ref{sec:fgap4}), and the triple
verification (V1, V2, V3) of \S\ref{sec:verification}.

\paragraph{What the present paper does NOT establish.}
The paper does NOT establish (a) an unconditional Pillar-4 closure
free of $\mathcal{H}_{\rm atomic}^{\rm final}$; (b) the numerical
F-GAP4 verdict itself, which awaits the Task~\#156 HYM-solver output
at the 2026-05-14 hard gate; (c) the exact $c_2(E)$ value on
$\Sigma_0$ at the convention-fixed level (Math305 \S3 numerical
caveat, $Q$-2026-05-29-Math305-c2-Sigma0-Numerical follow-up);
(d) the Bogomolov-type Proposition 4.3.1 of Math288 in fully
rigorous Hodge-theoretic form
($Q$-2026-05-01-Math288-Proposition-4-3-1-Rigorous-Proof
follow-up); (e) Pillar-4 closure on the original $\mathbb{CP}^2$
base, which Math174 falsified at $c_2(E) = -40$ and which motivated
the transfer to $\Sigma_0$ in the first place.

\paragraph{Implications for downstream pillars and Stage-2.}
The Pillar-4 atomic T6 PROVED CONDITIONAL tier sustains downstream
results: GAP-2 BRST closure (Math280) and GAP-3 anomaly cancellation
(Math281) are unconditional upgrades post-Pillar-4. The Stage-2
composite tier remains T3 pending the joint event $C_1 \wedge C_2
\wedge C_3$ on the F-GAP1 / F-GAP4 / F-Pillar6 verdict-shell
schedule per Math307 + Math310 + Math311; the present Pillar-4
atomic does not by itself promote Stage-2.

\paragraph{Historical $\mathbb{CP}^2$-route auxiliary record.}
The historical $\mathbb{CP}^2$-based version of the Pillar-4
programme (Math162 initial \v{C}ech construction, Math174 explicit
$c_2 = -40$ Chern-class calculation, Math191--192 principal-connection
Chern data) is preserved as a background algebro-geometric context
in the companion top-impact paper TI-1 (corrected
spin-c index formula on $\mathbb{CP}^2$). The Math270 topological
certificate transfer theorem provides the rigorous bridge from the
$\mathbb{CP}^2$ historical programme to the canonical $\Sigma_0$
present programme; the present paper operates on $\Sigma_0$
throughout.

% =====================================================================
\begin{acknowledgments}
The author thanks the Anthropic Claude Opus 4.7 collaborator for the
synthesis of the Pillar-4 canonical content from Math162--170,
Math191--229, Math248--281, Math288, Math300, Math302, Math305,
Math311 into the present revision (Rev 3, 2026-05-03), and for the
Math314 family hostile-referee audit cycle that drove the prior
Rev 1/Rev 2 corrections. This work is part of TECT, canonical
archive at \url{https://github.com/TECT-OpenPhysics/TECT}.
\end{acknowledgments}

% =====================================================================
\bibliographystyle{apsrev4-2}
\bibliography{../_shared/tect-references}
% =====================================================================

\end{document}
